\section{Translation Memory Redesign}
\label{sec:tmdesign}

A Translation Memory (TM) is a store of translations the pipeline has already made. Before sending a string to the model,
the pipeline looks it up; a \emph{hit} is reused instantly, and a \emph{miss} is translated and saved for next time. The
previous semester's memory worked, and measured a seeded 0\% versus 50\% cache. This chapter rebuilds it for the way the
pipeline is actually used: on every commit of an evolving application.

\subsection{Problems with the Previous Memory}
\begin{table}[H]
\centering
\caption[Problems of the previous memory]{Problems of the previous Translation Memory.}
\label{tab:tmflaws}
\begin{tabularx}{\linewidth}{@{}l L@{}}
\toprule
\textbf{Problem} & \textbf{Consequence} \\
\midrule
Emptied before every run & each commit retranslated the whole application, although a commit changes only a few strings \\
No check of which model made an entry & after a model or setting change, older translations would still be reused \\
One query per 1,000 strings per file & every file asked the database again for its own strings \\
Exact matches only & ``8 Cards'' could not reuse ``20 Cards''; 526 KDE4 strings (5.2\%) have such a near-twin \\
\bottomrule
\end{tabularx}
\end{table}

\subsection{Design}
\paragraph{Persistent, engine-labelled memory.} The memory is kept between runs. Each entry records its \emph{engine
version}: the model and every setting that changes the output (beam size, length cap, placeholder style, handling of
\code{<unk>}). Lookups match the current engine version, so a change of model or settings automatically stops older entries
from being reused, without deleting anything. The table key is (language, engine version, text hash), with the language
first so that one language's entries are stored together.

\paragraph{One load per language per run.} At the first file of a language, the pipeline loads that language's entries
for the current engine in a single query and keeps them in memory. Every later check is an in-process dictionary lookup,
and translations made during the run are added as they are saved, so later files reuse them without another query.

\paragraph{Near-twin reuse.} When the exact text is not found, the pipeline looks for a saved \emph{model} translation of a
string that differs only in numbers, placeholders or trailing punctuation, and adapts it:

\begin{center}
\small
\begin{tabular}{@{}l l l@{}}
\toprule
\textbf{Saved} & \textbf{New string} & \textbf{Adapted translation} \\
\midrule
\code{20 Cards} $\rightarrow$ \code{20 Karten} & \code{8 Cards} & \code{8 Karten} \\
\code{Copy \%1 to \%2} $\rightarrow$ \code{Nach \%2 kopieren: \%1} & \code{Copy \%3 to \%4} & \code{Nach \%4 kopieren: \%3} \\
\code{Border} $\rightarrow$ \code{Rahmen} & \code{Border:} & \code{Rahmen:} \\
\bottomrule
\end{tabular}
\end{center}

Reuse is refused, and the string goes to the model as usual, whenever it cannot be done cleanly: a value does not appear
exactly once in the saved translation (for example, the model wrote the number in local digits), the English repeats a
value (``1 of 1''), or the saved entry was itself adapted. Twins inside one file are sent to the model once and the rest
are adapted from that translation.

\subsection{Quality Check of Near-Twin Reuse}
Reuse only helps if adapted translations are as good as the model's own. Before enabling it, every twin pair in the full
KDE4 set (440 groups, 968 strings, 528 reusable pairs) was translated in all twelve languages both ways: freshly by the
production model and settings, and adapted from its twin. Both were scored against the references.

A first version \emph{guessed} how each language writes trailing punctuation and scored lower than fresh translation for
Japanese. Counting the endings of every KDE4 reference showed that only Chinese consistently uses the full-width colon,
that Japanese uses the ASCII colon slightly more often (749 to 657), that ``\ldots'' is kept in every language, and that
Hindi, Bengali and Punjabi end sentences with the danda while Marathi keeps the full stop. The
punctuation rules now follow these counts. Table~\ref{tab:twins} reports the final rules.

\begin{table}[H]
\centering
\caption[Near-twin reuse against fresh translation]{Near-twin reuse against fresh model translation (chrF on the reused strings, all 528 pairs per language).}
\label{tab:twins}
\begin{tabular}{@{}l S[table-format=3.0] S[table-format=2.0] S[table-format=2.2] S[table-format=2.2]@{}}
\toprule
\textbf{Language} & {\textbf{Reused}} & {\textbf{Refused}} & {\textbf{chrF fresh}} & {\textbf{chrF reused}} \\
\midrule
Hindi & 506 & 22 & 67.12 & 70.62 \\
Arabic & 488 & 40 & 55.51 & 56.08 \\
Bengali & 454 & 74 & 64.10 & 65.34 \\
German & 526 & 2 & 66.12 & 66.83 \\
Spanish & 527 & 1 & 66.82 & 68.04 \\
French & 510 & 18 & 67.54 & 68.74 \\
Japanese & 525 & 3 & 50.62 & 51.14 \\
Marathi & 492 & 36 & 57.62 & 59.55 \\
Punjabi & 527 & 1 & 60.24 & 62.60 \\
Tamil & 526 & 2 & 66.62 & 68.28 \\
Telugu & 526 & 2 & 66.98 & 70.50 \\
Chinese & 477 & 51 & 33.78 & 43.82 \\
\midrule
\textbf{Mean of 12} & & & \bfseries 60.26 & \bfseries 62.63 \\
\bottomrule
\end{tabular}
\end{table}

\begin{figure}[H]
\centering
\begin{tikzpicture}
\begin{axis}[reportbars, width=\linewidth, height=5.4cm, bar width=9pt,
  symbolic x coords={hi,ar,bn,de,es,fr,ja,mr,pa,ta,te,zh}, xtick=data,
  ylabel={chrF}, ymin=25, ymax=78, enlarge x limits=0.05]
\addplot[fill=cprev!75, draw=cprev] coordinates {(hi,67.12) (ar,55.51) (bn,64.10) (de,66.12) (es,66.82) (fr,67.54) (ja,50.62) (mr,57.62) (pa,60.24) (ta,66.62) (te,66.98) (zh,33.78)};
\addplot[fill=cnew!80, draw=cnew] coordinates {(hi,70.62) (ar,56.08) (bn,65.34) (de,66.83) (es,68.04) (fr,68.74) (ja,51.14) (mr,59.55) (pa,62.60) (ta,68.28) (te,70.50) (zh,43.82)};
\legend{Fresh model translation, Adapted from a twin}
\end{axis}
\end{tikzpicture}
\caption[Near-twin reuse per language]{Near-twin reuse scores at least as well as fresh translation in every language.}
\label{fig:twins}
\end{figure}

\paragraph{Finding.} By chrF, adapted translations score \textbf{equal or higher in all twelve languages} (mean 62.63
against 60.26), with placeholders kept 100\% in both. chrF rewards consistency: the model translates each string in
isolation and may translate ``Points'' one way in ``42 Points'' and another in ``78 Points'', while reuse applies the same
translation to every twin. A COMET check on the strings where reuse actually changed the output
(Chapter~\ref{sec:tmexp}) shows \textbf{equal} quality rather than higher, so the benefit of reuse is saved model work
and a more consistent application, at unchanged quality. Refusals
(8\%) were mostly strings whose numbers the model wrote in local digits; those strings were simply translated as before.
Near-twin reuse is enabled by default.
